% c15 原书 p.181–194（PDF p.188–201）
\begingroup
\setlength{\emergencystretch}{2.5em}
\subsection{仿射概形与射影概形的上同调}
\label{sec:2.4}

在计算仿射概形与射影概形的上同调之前，我们先研究 Koszul 复形的上同调。

设 $A$ 是环，$f_1,\dots,f_r\in A$。如下定义 \emph{Koszul 复形}\idx{koszulfuxing}{Koszul 复形} $K_\cdot(f_1,\dots,f_r)$：$K_1$ 是秩为 $r$ 的自由 $A$-模 $A^r$，且对任意 $p\ge 0$ 有 $K_p=\Lambda^p K_1$。对每个 $i\in\{1,\dots,r\}$，设 $e_i$ 是 $K_1=A^r$ 中第 $i$ 个分量为 $1$、其余分量为 $0$ 的元素。则 $e_i$（$i=1,\dots,r$）构成 $K_1$ 的一组基，而 $e_{i_1}\wedge\cdots\wedge e_{i_p}$（$i_1<\cdots<i_p$）构成 $K_p=\Lambda^p K_1$ 的一组基。定义 $\partial_p\colon K_p\to K_{p-1}$ 为用 $f_1e_1^{*}+\cdots+f_re_r^{*}$ 作缩并，其中 $\{e_1^{*},\dots,e_r^{*}\}$ 是 $\Hom_A(A^r,A)$ 中与 $\{e_1,\dots,e_r\}$ 对偶的基。我们有
\[
\begin{aligned}
\partial(e_{i_1}\wedge\cdots\wedge e_{i_p})
&=(f_1e_1^{*}+\cdots+f_re_r^{*})\lrcorner\,(e_{i_1}\wedge\cdots\wedge e_{i_p})\\
&=\sum_{k=1}^{p}(-1)^{k-1}f_{i_k}e_{i_1}\wedge\cdots\wedge\hat{e}_{i_k}\wedge\cdots\wedge e_{i_p}.
\end{aligned}
\]

对任意 $A$-模 $M$，定义
\[
\begin{aligned}
K_\cdot(f_1,\dots,f_r;M)&=K_\cdot(f_1,\dots,f_r)\tensor_A M,\\
K^{\cdot}(f_1,\dots,f_r;M)&=\Hom_A\bigl(K_\cdot(f_1,\dots,f_r),M\bigr),\\
H_p(f_1,\dots,f_p;M)&=H_p\bigl(K_\cdot(f_1,\dots,f_r;M)\bigr),\\
H^p(f_1,\dots,f_p;M)&=H^p\bigl(K^{\cdot}(f_1,\dots,f_r;M)\bigr).
\end{aligned}
\]

\begin{tnote}
上式中 $H_p$ 与 $H^p$ 的第二个宗量在原书 p.182 印作 $f_p$，但等号右端是复形 $K_\cdot(f_1,\dots,f_r;M)$ 的同调，依本段开头及 2.4.1，此处应为 $f_r$。[依据：原书 p.182 本段首行与命题 2.4.1]
\end{tnote}

$K^p(f_1,\dots,f_r;M)=\Hom_A(K_p,M)$ 中的任一元素 $g$ 完全由
\[
g_{i_1\dots i_p}=g(e_{i_1}\wedge\cdots\wedge e_{i_p})
\]
确定。我们有
\[
\begin{aligned}
(dg)_{i_1\dots i_{p+1}}&=(dg)(e_{i_1}\wedge\cdots\wedge e_{i_{p+1}})
=g\bigl(\partial(e_{i_1}\wedge\cdots\wedge e_{i_{p+1}})\bigr)\\
&=\sum_{k=1}^{p+1}(-1)^{k-1}f_{i_k}g_{i_1\dots\hat{i_k}\dots i_{p+1}}.
\end{aligned}
\]
我们有
\[
K^p(f_1,\dots,f_r;M)=\Hom_A(K_p,M)\cong\Lambda^p(A^r)^{*}\tensor_A M.
\]
我们可以写
\[
g=\sum_{i_1<\cdots<i_p}e_{i_1}^{*}\wedge\cdots\wedge e_{i_p}^{*}\tensor g_{i_1\dots i_p}.
\]
我们有
\[
dg=(f_1e_1^{*}+\cdots+f_re_r^{*})\wedge g.
\]
事实上，我们有
\[
\begin{aligned}
&(f_1e_1^{*}+\cdots+f_re_r^{*})\wedge\Bigl(\sum_{i_1<\cdots<i_p}e_{i_1}^{*}\wedge\cdots\wedge e_{i_p}^{*}\tensor g_{i_1\dots i_p}\Bigr)\\
&\qquad=\sum_{i_1<\cdots<i_p}\ \sum_{j\in\{1,\dots,r\}-\{i_1,\dots,i_p\}}
f_je_j^{*}\wedge e_{i_1}^{*}\wedge\cdots\wedge e_{i_p}^{*}\tensor g_{i_1\dots i_p}\\
&\qquad=\sum_{i_1<\cdots<i_{p+1}}e_{i_1}^{*}\wedge\cdots\wedge e_{i_{p+1}}^{*}
\tensor\Bigl(\sum_{k=1}^{p+1}(-1)^{k-1}f_{i_k}g_{i_1\dots\hat{i_k}\dots i_{p+1}}\Bigr).
\end{aligned}
\]

\begin{proposition}\label{it:2.4.1}
对每个 $p$，有 $H_p(f_1,\dots,f_r;M)\cong H^{r-p}(f_1,\dots,f_r;M)$。
\end{proposition}

\begin{proof}
考虑 $A$-模同态
\[
* \colon K_p(f_1,\dots,f_r;M)\to K^{r-p}(f_1,\dots,f_r;M),
\]
\[
*(e_{i_1}\wedge\cdots\wedge e_{i_p})=\operatorname{sgn}(i_1\dots i_pj_1\dots j_{r-p})\,e_{j_1}^{*}\wedge\cdots\wedge e_{j_{r-p}}^{*},
\]
其中 $\{i_1,\dots,i_p,j_1,\dots,j_{r-p}\}=\{1,\dots,r\}$，而 $\operatorname{sgn}(i_1\dots i_pj_1\dots j_{r-p})$ 是置换
\[
\begin{pmatrix}
1 & \dots & p & p+1 & \dots & r\\
i_1 & \dots & i_p & j_1 & \dots & j_{r-p}
\end{pmatrix}
\]
的符号。

我们断言，在 $K_p(f_1,\dots,f_r;M)$ 上有
\[
*\partial=(-1)^{p-1}d*.
\]
我们的断言随即成立。事实上，有
\[
\begin{aligned}
&*\partial(e_{i_1}\wedge\cdots\wedge e_{i_p})\\
&\quad=*(f_1e_1^{*}+\cdots+f_re_r^{*})\lrcorner\,(e_{i_1}\wedge\cdots\wedge e_{i_p})\\
&\quad=*\Bigl(\sum_{k=1}^{p}(-1)^{k-1}f_{i_k}e_{i_1}\wedge\cdots\wedge\hat{e}_{i_k}\wedge\cdots\wedge e_{i_p}\Bigr)\\
&\quad=\sum_{k=1}^{p}(-1)^{k-1}f_{i_k}
\operatorname{sgn}(i_1\dots\hat{i_k}\dots i_pi_kj_1\dots j_{r-p})e_{i_k}^{*}\wedge e_{j_1}^{*}\wedge\cdots\wedge e_{j_{r-p}}^{*}\\
&\quad=(-1)^{p-1}\operatorname{sgn}(i_1\dots i_pj_1\dots j_{r-p})
\sum_{k=1}^{p}f_{i_k}e_{i_k}^{*}\wedge e_{j_1}^{*}\wedge\cdots\wedge e_{j_{r-p}}^{*},
\end{aligned}
\]
\[
\begin{aligned}
&d*(e_{i_1}\wedge\cdots\wedge e_{i_p})\\
&\quad=\operatorname{sgn}(i_1\dots i_pj_1\dots j_{r-p})d(e_{j_1}^{*}\wedge\cdots\wedge e_{j_{r-p}}^{*})\\
&\quad=\operatorname{sgn}(i_1\dots i_pj_1\dots j_{r-p})(f_1e_1^{*}+\cdots+f_re_r^{*})\wedge(e_{j_1}^{*}\wedge\cdots\wedge e_{j_{r-p}}^{*})\\
&\quad=\operatorname{sgn}(i_1\dots i_pj_1\dots j_{r-p})\sum_{k=1}^{p}f_{i_k}e_{i_k}^{*}\wedge e_{j_1}^{*}\wedge\cdots\wedge e_{j_{r-p}}^{*}.
\end{aligned}
\]
\end{proof}

\begin{proposition}\label{it:2.4.2}
设 $f_1,\dots,f_r$ 生成 $A$ 的单位理想。则对每个 $A$-模 $M$，有 $H_p(f_1,\dots,f_r;M)=0$ 且 $H^p(f_1,\dots,f_r;M)=0$ 对所有 $p$ 成立。
\end{proposition}

\begin{proof}
由 \ref{it:2.4.1}，只需证明 $H^p(f_1,\dots,f_r;M)=0$。取 $a_1,\dots,a_r\in A$ 使得 $a_1f_1+\cdots+a_rf_r=1$。定义
\[
H\colon K^p(f_1,\dots,f_r;M)\to K^{p-1}(f_1,\dots,f_r;M)
\]
为用 $a_1e_1+\cdots+a_re_r$ 作缩并。对任意 $g\in K^p(f_1,\dots,f_r;M)$，有
\[
\begin{aligned}
(dH+Hd)g
&=\Bigl(\sum_{k=1}^{p}f_ke_k^{*}\Bigr)\wedge\Bigl(\Bigl(\sum_{l=1}^{p}a_le_l\Bigr)\lrcorner\, g\Bigr)
+\Bigl(\sum_{l=1}^{p}a_le_l\Bigr)\lrcorner\Bigl(\Bigl(\sum_{k=1}^{p}f_ke_k^{*}\Bigr)\wedge g\Bigr)\\
&=\sum_{k,l=1}^{p}a_lf_ke_k^{*}\wedge(e_l\lrcorner\, g)+\sum_{k,l=1}^{p}a_lf_k(e_l\lrcorner\, e_k^{*})g
-\sum_{k,l=1}^{p}a_lf_ke_k^{*}\wedge(e_l\lrcorner\, g)\\
&=\sum_{k=1}^{p}a_kf_kg=g.
\end{aligned}
\]
所以 $dH+Hd=\id$，即 $K^{\cdot}(f_1,\dots,f_r;M)$ 上的恒等态射同伦于 $0$。这说明 $H^p(f_1,\dots,f_r;M)=0$ 对所有 $p$ 成立。
\end{proof}

\begin{tnote}
2.4.2 证明中的求和上限在原书 p.183 都印作 $p$（如 $\sum_{k=1}^{p}f_ke_k^{*}$、$\sum_{k,l=1}^{p}$）；但缩并所用的元素是 $a_1e_1+\cdots+a_re_r$，且 $g$ 的次数与求和范围无关，故诸求和的上限应为 $r$。[依据：原书 p.183 同页 $H$ 的定义式 $a_1e_1+\cdots+a_re_r$]
\end{tnote}

\begin{proposition}\label{it:2.4.3}
设 $f_1,\dots,f_r\in A$，$M$ 是 $A$-模。假设对每个 $1\le k\le r$，同态
\[
M\Big/\sum_{i=1}^{k-1}f_iM\to M\Big/\sum_{i=1}^{k-1}f_iM,\qquad z\mapsto f_kz
\]
是单射。则 $H_p(f_1,\dots,f_r;M)=0$ 对每个 $p\ne 0$ 成立，且 $H^p(f_1,\dots,f_r;M)=0$ 对每个 $p\ne r$ 成立。
\end{proposition}

\begin{proof}
由 \ref{it:2.4.1}，只需证明 $H_p(f_1,\dots,f_r;M)=0$ 对每个 $p\ne 0$ 成立。我们对 $r$ 归纳证明这一点。

当 $r=1$ 时，$K_\cdot(f_1;M)$ 等同于复形
\[
0\to M\xrightarrow{\;f_1\;}M\to 0,
\]
其中态射 $f_1$ 是乘以 $f_1$。由于 $f_1\colon M\to M$ 是单射，有 $H_p(f_1;M)=0$ 对每个 $p\ne 0$ 成立。

设已证 $H_p(f_1,\dots,f_{r-1};M)=0$ 对每个 $p\ne 0$ 成立。注意 $K_\cdot(f_1,\dots,f_{r-1};M)$ 可以看作 $K_\cdot(f_1,\dots,f_r;M)$ 的子复形。此外，有同构
\[
K_{p-1}(f_1,\dots,f_{r-1};M)\to K_p(f_1,\dots,f_r;M)/K_p(f_1,\dots,f_{r-1};M),
\]
\[
e_{i_1}\wedge\cdots\wedge e_{i_{p-1}}\tensor z\mapsto
e_{i_1}\wedge\cdots\wedge e_{i_{p-1}}\wedge e_r\tensor z
\]
\[
(z\in M,\quad 1\le i_1<\cdots<i_{p-1}\le r-1),
\]
且这些同构与 $\partial$ 交换。于是有复形的正合序列
\[
0\to K_\cdot(f_1,\dots,f_{r-1};M)\to K_\cdot(f_1,\dots,f_r;M)\to K_\cdot(f_1,\dots,f_{r-1};M)[-1]\to 0,
\]
其中复形 $K_\cdot(f_1,\dots,f_{r-1};M)[-1]$ 由
\[
\bigl(K_\cdot(f_1,\dots,f_{r-1};M)[-1]\bigr)_p=K_{p-1}(f_1,\dots,f_{r-1};M)
\]
定义。考虑与此短正合序列相联系的上同调群的长正合序列
\[
\cdots\to H_p(f_1,\dots,f_{r-1};M)\to H_p(f_1,\dots,f_r;M)\to H_{p-1}(f_1,\dots,f_{r-1};M)\to\cdots.
\]
由归纳假设，有
\[
H_p(f_1,\dots,f_{r-1};M)=H_{p-1}(f_1,\dots,f_{r-1};M)=0
\]
对每个 $p\ge 2$ 成立。所以 $H_p(f_1,\dots,f_r;M)=0$ 对每个 $p\ge 2$ 成立。此外，上述长正合序列的一部分是
\[
0\to H_1(f_1,\dots,f_r;M)\to H_0(f_1,\dots,f_{r-1};M)\to H_0(f_1,\dots,f_{r-1};M).
\]
可以验证，最后一个箭头可以等同于
\[
M\Big/\sum_{i=1}^{r-1}f_rM\xrightarrow{\;f_r\;}M\Big/\sum_{i=1}^{r-1}f_rM.
\]
由于它是单射，有 $H_1(f_1,\dots,f_r;M)=0$。
\end{proof}

\begin{tnote}
2.4.3 证明末尾的映射在原书 p.184 印作 $M/\sum_{i=1}^{r-1}f_rM\to M/\sum_{i=1}^{r-1}f_rM$；但按命题条件，$f_r$ 应作用于商模 $M/\sum_{i=1}^{r-1}f_iM$，故两处求和项中的 $f_r$ 疑为 $f_i$。[依据：原书 p.184 命题 2.4.3 的条件式]
\end{tnote}

\begin{proposition}\label{it:2.4.4}
设 $X=\Spec A$ 是仿射概形，$\cF$ 是拟凝聚 $\cO_X$-模。则对每个 $p\ge 1$，有 $\check{H}^p(X,\cF)=0$。
\end{proposition}

\begin{proof}
对每个 $A$-模 $M$ 与每个 $f\in A$，考虑有向系统 $(M^{(n)},f_{mn})_{n\in\bN}$，其中 $M^{(n)}=M$ 对每个 $n\in\bN$ 成立，且对每对 $m\le n$，
\[
f_{mn}\colon M^{(m)}\to M^{(n)},\quad x\mapsto f^{n-m}x.
\]
同态族
\[
f_n\colon M^{(n)}\to M_f,\quad x\mapsto\frac{x}{f^n}
\]
诱导同构
\[
\dlim_n M^{(n)}\cong M_f.
\]

设 $\cF\cong M^{\sim}$。取 $f_1,\dots,f_r\in A$ 使得 $\fU=\{D(f_1),\dots,D(f_r)\}$ 是 $X$ 的开覆盖。交错 \v{C}ech 复形由
\[
\begin{aligned}
C^p(\fU,\cF)
&=\prod_{1\le i_0<\cdots<i_p\le r}\cF\bigl(D(f_{i_0})\cap\cdots\cap D(f_{i_p})\bigr)\\
&=\prod_{1\le i_0<\cdots<i_p\le r}M_{f_{i_0}\cdots f_{i_p}}
=\dlim_n\prod_{1\le i_0<\cdots<i_p\le r}M^{(n)}_{i_0\dots i_p}
\end{aligned}
\]
给出，其中有向系统 $\bigl(\prod_{1\le i_0<\cdots<i_p\le r}M^{(n)}_{i_0\dots i_p}\bigr)_{n\in\bN}$ 由
\[
\prod_{1\le i_0<\cdots<i_p\le r}M^{(n)}_{i_0\dots i_p}
=\prod_{1\le i_0<\cdots<i_p\le r}M
\]
对每个 $n$ 定义，并且对每对 $m\le n$，
\[
\prod_{1\le i_0<\cdots<i_p\le r}M^{(m)}_{i_0\dots i_p}
\to\prod_{1\le i_0<\cdots<i_p\le r}M^{(n)}_{i_0\dots i_p},\quad
(x_{i_0\dots i_p})\mapsto\bigl((f_{i_0}\cdots f_{i_p})^{n-m}x_{i_0\dots i_p}\bigr).
\]
令 $C^p_n(M)=\prod_{1\le i_0<\cdots<i_p\le r}M^{(n)}_{i_0\dots i_p}$，并定义
\[
d\colon C^p_n(M)\to C^{p+1}_n(M),\quad x=(x_{i_0\dots i_p})\mapsto dx=\bigl((dx)_{i_0\dots i_{p+1}}\bigr),
\]
其中
\[
(dx)_{i_0\dots i_{p+1}}=\sum_{k=0}^{p+1}(-1)^kf_{i_k}^nx_{i_0\dots\hat{i_k}\dots i_{p+1}}.
\]
则每个 $C_n(M)$ 是复形，且
\[
C^{\cdot}(\fU,\cF)\cong\dlim_n C_n(M).
\]
所以对每个 $p$，有
\[
\check{H}^p(\fU,\cF)\cong\dlim_n H^p(C_n^{\cdot}(M)).
\]
我们将证明 $H^p(C_n^{\cdot}(M))=0$ 对每个 $n$ 与每个 $p\ge 1$ 成立。这蕴含 $\check{H}^p(\fU,\cF)=0$ 对每个 $p\ge 1$ 成立。由于 $X$ 的任一开覆盖都有由形如 $D(f)$（$f\in A$）的开子集组成的加细，有 $\check{H}^p(X,\cF)=0$ 对每个 $p\ge 1$ 成立。

由于 $D(f_1^n),\dots,D(f_r^n)$ 覆盖 $X=\Spec A$，$f_1^n,\dots,f_r^n$ 生成 $A$ 的单位理想。注意对每个 $p\ge 1$，序列
\[
C_n^{p-1}(M)\xrightarrow{\;d\;}C_n^p(M)\xrightarrow{\;d\;}C_n^{p+1}(M)
\]
可以等同于序列
\[
K^p(f_1^n,\dots,f_r^n;M)\xrightarrow{\;d\;}K^{p+1}(f_1^n,\dots,f_r^n;M)
\xrightarrow{\;d\;}K^{p+2}(f_1^n,\dots,f_r^n;M).
\]
所以我们有
\[
H^p(C_n^{\cdot}(M))\cong H^{p+1}(f_1^n,\dots,f_r^n;M).
\]
我们的断言随即由 \ref{it:2.4.2} 得到。
\end{proof}

\begin{corollary}\label{it:2.4.5}
设 $X$ 是分离概形，$\cF$ 是拟凝聚 $\cO_X$-模。则对 $X$ 的每个由仿射开子概形组成的覆盖 $\fU$，有 $\check{H}^p(\fU,\cF)\cong H^p(X,\cF)$ 对每个 $p$ 成立。此外有 $\check{H}^p(X,\cF)\cong H^p(X,\cF)$ 对每个 $p$ 成立。
\end{corollary}

\begin{proof}
设 $\fB$ 是 $X$ 的仿射开子概形组成的族。则对 $\fB$ 中任意 $U_1,\dots,U_k$，由 \ref{it:1.3.25}(i)，$U_1\cap\cdots\cap U_k$ 也是仿射的，从而由 \ref{it:2.4.4} 有 $\check{H}^q(U_1\cap\cdots\cap U_k,\cF)=0$ 对每个 $q\ge 1$ 成立。于是推论由 \ref{it:2.3.11} 得到。
\end{proof}

\begin{theorem}\label{it:2.4.6}
设 $X$ 是仿射概形，$\cF$ 是拟凝聚 $\cO_X$-模。则对每个 $p\ge 1$，有 $H^p(X,\cF)=0$。
\end{theorem}

\begin{proof}
由 \ref{it:2.4.5}，有 $H^p(X,\cF)=\check{H}^p(X,\cF)$ 对每个 $p$ 成立；由 \ref{it:2.4.4}，有 $\check{H}^p(X,\cF)=0$ 对每个 $p\ge 1$ 成立。断言得证。
\end{proof}

\begin{corollary}\label{it:2.4.7}
设
\[
0\to\cF\to\cG\to\cH\to 0
\]
是概形 $X$ 上 $\cO_X$-模的正合序列。若 $\cF$ 与 $\cH$ 是拟凝聚的，则 $\cG$ 也是拟凝聚的。
\end{corollary}

\begin{proof}
问题在局部，故可设 $X$ 是仿射的。由于 $\cF$ 是拟凝聚的，由 \ref{it:2.4.6} 有 $H^1(X,\cF)=0$。于是与上述短正合序列相伴的层上同调群的长正合序列给出正合序列
\[
0\to\Gamma(X,\cF)\to\Gamma(X,\cG)\to\Gamma(X,\cH)\to 0.
\]
考虑交换图
\begin{center}
\begin{tikzcd}
0 \arrow[r] & \Gamma(X,\cF)^{\sim} \arrow[r] \arrow[d] & \Gamma(X,\cG)^{\sim} \arrow[r] \arrow[d] & \Gamma(X,\cH)^{\sim} \arrow[r] \arrow[d] & 0\\
0 \arrow[r] & \cF \arrow[r] & \cG \arrow[r] & \cH \arrow[r] & 0.
\end{tikzcd}
\end{center}
第一与第三个竖直箭头是同构，因为 $\cF$ 与 $\cH$ 是拟凝聚的。所以中间的竖直箭头也是同构，从而 $\cG$ 是拟凝聚的。
\end{proof}

\begin{corollary}\label{it:2.4.8}
设 $f\colon X\to Y$ 是仿射态射，$\cF$ 是拟凝聚 $\cO_X$-模。则对每个 $p\ge 1$，有 $R^pf_*\cF=0$。
\end{corollary}

\begin{proof}
对 $Y$ 的每个仿射开子概形 $V$，有
\[
H^p(f^{-1}(V),\cF)=0
\]
对每个 $p\ge 1$ 成立，因为 $f^{-1}(V)$ 是仿射的。断言随即由 \ref{it:2.1.9} 以及仿射开子集构成 $Y$ 的拓扑的一组基这一事实得到。
\end{proof}

设 $\fU=\{U_i\}_{i\in I}$ 是 $X$ 的开覆盖。对 $\fU$ 中任意开子集 $U_{i_0},\dots,U_{i_n}$，用 $j_{i_0\dots i_n}\colon U_{i_0\dots i_n}\hookrightarrow X$ 记包含态射。如下定义\emph{层化的 \v{C}ech 复形}\idx{cengtadecechfuxing}{层化的 \v{C}ech 复形} $C^{\cdot}(\fU,\cF)$：在 $I$ 上固定一个全序。对每个 $n\ge 0$，令
\[
C^n(\fU,\cF)=\prod_{i_0<\cdots<i_n}j_{i_0\dots i_n*}\cF.
\]
定义
\[
d\colon C^n(\fU,\cF)\to C^{n+1}(\fU,\cF)
\]
如下：对 $X$ 的每个开子集 $V$ 与 $C^n(\fU,\cF)(V)=\prod_{i_0<\cdots<i_n}\cF(U_{i_0\dots i_n}\cap V)$ 中的每个截面 $s=(s_{i_0\dots i_n})$，定义 $C^{n+1}(\fU,\cF)(V)=\prod_{i_0<\cdots<i_{n+1}}\cF(U_{i_0\dots i_{n+1}}\cap V)$ 中的截面 $ds=((ds)_{i_0\dots i_{n+1}})$ 为
\[
(ds)_{i_0\dots i_{n+1}}=\sum_{k=0}^{n+1}(-1)^ks_{i_0\dots\hat{i_k}\dots i_{n+1}}\big|_{U_{i_0\dots i_{n+1}}\cap V}.
\]
注意 $C^{\cdot}(\fU,\cF)$ 的第 $n$ 个上同调层是与预层 $V\mapsto\check{H}^n(\fU\cap V,\cF)$ 相伴的层，其中 $\fU\cap V$ 是 $V$ 的开覆盖 $\{U_i\cap V\}_{i\in I}$。

\begin{proposition}\label{it:2.4.9}
设 $f\colon X\to Y$ 是概形的拟紧且分离的态射。则对每个拟凝聚 $\cO_X$-模 $\cF$ 与每个 $n$，$\cO_Y$-模 $R^nf_*\cF$ 是拟凝聚的。此外，对 $Y$ 的每个仿射开子概形 $V$，有典范同构
\[
H^n(f^{-1}(V),\cF)\cong\Gamma(V,R^nf_*\cF).
\]
若 $Y=\Spec A$ 是仿射的，则
\[
R^nf_*\cF\cong H^n(X,\cF)^{\sim}.
\]
\end{proposition}

\begin{proof}
可设 $Y$ 是仿射的。设 $\fU$ 是 $X$ 的由有限多个仿射开子概形组成的覆盖，$C^{\cdot}(\fU,\cF)$ 是层化的 \v{C}ech 复形。考虑 $Y$ 上层的复形 $f_*C^{\cdot}(\fU,\cF)$。它的第 $n$ 个上同调层 $H^n(f_*C^{\cdot}(\fU,\cF))$ 是与预层
\[
V\mapsto\check{H}^n(\fU\cap f^{-1}(V),\cF)
\]
相伴的层。

当 $V$ 是 $Y$ 的仿射开子概形时，$\fU\cap f^{-1}(V)$ 是 $f^{-1}(V)$ 的由仿射开子概形组成的覆盖。所以由 \ref{it:2.4.5}，有
\[
\check{H}^n(\fU\cap f^{-1}(V),\cF)\cong H^n(f^{-1}(V),\cF)
\]
对 $Y$ 的每个仿射开子概形 $V$ 成立。因此 $H^n(f_*C^{\cdot}(\fU,\cF))$ 也是与预层
\[
V\mapsto H^n(f^{-1}(V),\cF)
\]
相伴的层。

于是由 \ref{it:2.1.9}，有
\[
R^nf_*\cF\cong H^n(f_*C^{\cdot}(\fU,\cF)).
\]
由 \ref{it:1.4.9}(iii)，每个 $f_*C^n(\fU,\cF)$ 是拟凝聚的。所以 $R^nf_*\cF$ 也是拟凝聚的。此外，由于 $Y$ 是仿射的，有
\[
f_*C^{\cdot}(\fU,\cF)\cong\Gamma(Y,f_*C^{\cdot}(\fU,\cF))^{\sim}.
\]
我们有
\[
\Gamma(Y,f_*C^{\cdot}(\fU,\cF))=C^{\cdot}(\fU,\cF).
\]
于是有
\[
\begin{aligned}
R^nf_*\cF&\cong H^n(C^{\cdot}(\fU,\cF))^{\sim}\\
&\cong\check{H}^n(\fU,\cF)^{\sim}\\
&\cong H^n(X,\cF)^{\sim}.
\end{aligned}
\]
特别地，有
\[
\Gamma(Y,R^nf_*\cF)\cong H^n(X,\cF).
\]
\end{proof}

概形态射 $g\colon Y'\to Y$ 称为\emph{平坦态射}\idx{pingtande}{平坦态射}，若对每个 $y'\in Y'$，$\cO_{Y',y'}$ 是平坦的 $\cO_{Y,g(y')}$-代数。由 \ref{it:2.1.19}，这等价于说：对 $Y'$ 的每个仿射开子概形 $U'$ 以及 $Y$ 的每个满足 $g(U')\subset U$ 的仿射开子概形 $U$，$\cO_{Y'}(U')$ 是平坦的 $\cO_Y(U)$-代数。

\begin{proposition}\label{it:2.4.10}
设 $f\colon X\to Y$ 是拟紧且分离的态射，$g\colon Y'\to Y$ 是平坦态射。用下述笛卡尔图固定记号：
\begin{center}
\begin{tikzcd}
X\times_YY' \arrow[r,"g'"] \arrow[d,"f'"] & X \arrow[d,"f"]\\
Y' \arrow[r,"g"] & Y.
\end{tikzcd}
\end{center}
则对 $X$ 上每个拟凝聚 $\cO_X$-模 $\cF$ 与每个 $p$，有典范同构
\[
g^{*}R^pf_*\cF\cong R^pf'_*g'^{*}\cF.
\]
\end{proposition}

\begin{proof}
设 $\cI^{\cdot}$ 是 $\cF$ 的内射分解（在 $\cO_X$-模范畴中），$\cJ^{\cdot}$ 是 $g'^{*}\cF$ 的内射分解（在 $\cO_{X\times_YY'}$-模范畴中）。由于 $g$ 平坦，$g'$ 也平坦。所以函子 $g'^{*}$ 正合，从而有正合序列
\[
0\to g'^{*}\cF\to g'^{*}\cI^{\cdot}.
\]
由 \ref{it:2.1.1}(ii)，有在合同伦意义下唯一的态射 $g'^{*}\cI^{\cdot}\to\cJ^{\cdot}$，它延拓 $\id_{g'^{*}\cF}$。它诱导态射 $f'_*g'^{*}\cI^{\cdot}\to f'_*\cJ^{\cdot}$。另一方面，典范态射 $\cI^{\cdot}\to g'_*g'^{*}\cI^{\cdot}$ 诱导态射
\[
f_*\cI^{\cdot}\to f_*g'_*g'^{*}\cI^{\cdot}\cong g_*f'_*g'^{*}\cI^{\cdot}.
\]
由于 $g^{*}$ 是 $g_*$ 的左伴随，这个态射诱导态射 $g^{*}f_*\cI^{\cdot}\to f'_*g'^{*}\cI^{\cdot}$。再与态射 $f'_*g'^{*}\cI^{\cdot}\to f'_*\cJ^{\cdot}$ 复合，得到态射 $g^{*}f_*\cI^{\cdot}\to f'_*\cJ^{\cdot}$。由于 $g$ 平坦，$g^{*}$ 是正合函子，所以 $g^{*}f_*\cI^{\cdot}$ 的第 $p$ 个上同调层是 $g^{*}R^pf_*\cF$。于是由态射 $g^{*}f_*\cI^{\cdot}\to f'_*\cJ^{\cdot}$ 在上同调层上诱导的同态是
\[
g^{*}R^pf_*\cF\to R^pf'_*g'^{*}\cF.
\]
我们来证明它是同构。问题对 $Y$ 与 $Y'$ 是局部的。故可设 $Y=\Spec A$、$Y'=\Spec A'$ 都是仿射的。由 \ref{it:2.4.9}，有
\[
R^pf_*\cF\cong H^p(X,\cF)^{\sim},\qquad
R^pf'_*g'^{*}\cF\cong H^p(X\times_YY',g'^{*}\cF)^{\sim}.
\]
由 \ref{it:1.4.3}(ii)，我们于是有
\[
g^{*}R^pf_*\cF\cong\bigl(A'\tensor_AH^p(X,\cF)\bigr)^{\sim}.
\]
为证明本命题，只需证明典范同态
\[
A'\tensor_AH^p(X,\cF)\to H^p(X\times_YY',g'^{*}\cF)
\]
是同构。设 $\fU=\{U_i\}_{i\in I}$ 是 $X$ 的由有限多个仿射开子概形组成的覆盖。则 $\fU'=\{U_i\times_YY'\}$ 是 $X\times_YY'$ 的由仿射开子概形组成的覆盖。由 \ref{it:2.4.5}，有
\[
H^p(X,\cF)\cong\check{H}^p(\fU,\cF),\qquad
H^p(X\times_YY',g'^{*}\cF)\cong\check{H}^p(\fU',g'^{*}\cF).
\]
可以证明
\[
A'\tensor_AC^{\cdot}(\fU,\cF)\cong C^{\cdot}(\fU',g'^{*}\cF).
\]
在两边取第 $p$ 个上同调群，并利用 $A'$ 在 $A$ 上平坦这一事实，得到 $A'\tensor_AH^p(X,\cF)\cong H^p(X\times_YY',g'^{*}\cF)$。
\end{proof}

\begin{proposition}\label{it:2.4.11}
设 $A$ 是环，$S$ 是环 $A[x_0,\dots,x_r]$ 并赋以由多项式的次数定义的分次，$X=\Proj S$。
\begin{enumerate}[label=(\roman*),leftmargin=2.2em]
  \item 典范同态
  \[
  S\to\Gamma_*(\cO_X)=\bigoplus_{n=-\infty}^{\infty}H^0(X,\cO_X(n))
  \]
  是同构。
  \item $\bigoplus_{n=-\infty}^{\infty}H^r(X,\cO_X(n))$ 是以 $\dfrac{1}{x_0^{k_0}\cdots x_r^{k_r}}$（$k_i>0$）为基的自由 $A$-模，其中 $\dfrac{1}{x_0^{k_0}\cdots x_r^{k_r}}$ 位于 $H^r(X,\cO_X(-(k_0+\cdots+k_r)))$ 中。
  \item 对每个 $p\ne 0,r$，有 $\bigoplus_{n=-\infty}^{\infty}H^p(X,\cO_X(n))=0$。
\end{enumerate}
\end{proposition}

\begin{proof}
(i) 在 \ref{it:1.4.23} 中证明。我们来证 (ii) 与 (iii)。考虑 $X$ 的开覆盖 $\fU=\{D_+(x_i)\mid i=0,\dots,r\}$。由 \ref{it:2.4.5}，对每个 $p$，有
\[
\bigoplus_{n=-\infty}^{\infty}H^p(X,\cO_X(n))\cong
\bigoplus_{n=-\infty}^{\infty}\check{H}^p\bigl(\fU,\cO_X(n)\bigr).
\]
考虑交错 \v{C}ech 复形
\[
C^{\cdot}\Bigl(\fU,\bigoplus_{n=-\infty}^{\infty}\cO_X(n)\Bigr)
=\bigoplus_{n=-\infty}^{\infty}C^{\cdot}\bigl(\fU,\cO_X(n)\bigr).
\]
我们有
\[
\begin{aligned}
C^{\cdot p}\Bigl(\fU,\bigoplus_{n=-\infty}^{\infty}\cO_X(n)\Bigr)
&=\prod_{0\le i_0<\cdots<i_p\le r}\Bigl(\bigoplus_{n=-\infty}^{\infty}\cO_X(n)\Bigr)
\bigl(D_+(x_{i_0}\cdots x_{i_p})\bigr)\\
&=\prod_{0\le i_0<\cdots<i_p\le r}S_{x_{i_0}\cdots x_{i_p}}.
\end{aligned}
\]
注意
\[
C^{r-1}\Bigl(\fU,\bigoplus_{n=-\infty}^{\infty}\cO_X(n)\Bigr)\to
C^r\Bigl(\fU,\bigoplus_{n=-\infty}^{\infty}\cO_X(n)\Bigr)\to
C^{r+1}\Bigl(\fU,\bigoplus_{n=-\infty}^{\infty}\cO_X(n)\Bigr)
\]
可以等同于
\[
\prod_{0\le i_0<\cdots<i_{r-1}\le r}S_{x_{i_0}\cdots x_{i_{r-1}}}\to S_{x_0\cdots x_r}\to 0.
\]
所以 $\bigoplus_{n=-\infty}^{\infty}H^r(X,\cO_X(n))$ 是同态
\[
\prod_{k=0}^{r}S_{x_0\cdots\hat{x}_k\cdots x_r}\to S_{x_0\cdots x_r},\quad
(g_k)\mapsto\sum_{k=0}^{r}(-1)^kg_k
\]
的余核。注意 $S_{x_0\cdots x_r}$ 是以 $x_0^{k_0}\cdots x_r^{k_r}$（$k_i\in\bZ$）为基的自由 $A$-模，而上述同态的像是由那些至少有一个 $k_i\ge 0$ 的 $x_0^{k_0}\cdots x_r^{k_r}$ 生成的子模。所以 $\bigoplus_{n=-\infty}^{\infty}H^r(X,\cO_X(n))$ 是以 $\dfrac{1}{x_0^{k_0}\cdots x_r^{k_r}}$（$k_i>0$）为基的自由 $A$-模。这就证明了 (ii)。

另一方面，我们有
\[
C^{\cdot p}\Bigl(\fU,\bigoplus_{n=-\infty}^{\infty}\cO_X(n)\Bigr)
=\prod_{0\le i_0<\cdots<i_p\le r}S_{x_{i_0}\cdots x_{i_p}}
=\dlim_n\prod_{0\le i_0<\cdots<i_p\le r}S^{(n)}_{i_0\dots i_p},
\]
其中有向系统 $\bigl(\prod_{0\le i_0<\cdots<i_p\le r}S^{(n)}_{i_0\dots i_p}\bigr)_{n\in\bN}$ 由
\[
\prod_{0\le i_0<\cdots<i_p\le r}S^{(n)}_{i_0\dots i_p}
=\prod_{0\le i_0<\cdots<i_p\le r}S
\]
对每个 $n$ 定义，并且对每对 $m\le n$，
\[
\prod_{0\le i_0<\cdots<i_p\le r}S^{(m)}_{i_0\dots i_p}
\to\prod_{0\le i_0<\cdots<i_p\le r}S^{(n)}_{i_0\dots i_p},\quad
(s_{i_0\dots i_p})\mapsto\bigl((x_{i_0}\cdots x_{i_p})^{n-m}s_{i_0\dots i_p}\bigr).
\]
令 $C^p_n(S)=\prod_{0\le i_0<\cdots<i_p\le r}S^{(n)}_{i_0\dots i_p}$，并定义
\[
d\colon C^p_n(S)\to C^{p+1}_n(S),\quad s=(s_{i_0\dots i_p})\mapsto ds=\bigl((ds)_{i_0\dots i_{p+1}}\bigr),
\]
其中
\[
(ds)_{i_0\dots i_{p+1}}=\sum_{k=0}^{p+1}(-1)^kx_{i_k}^ns_{i_0\dots\hat{i_k}\dots i_{p+1}}.
\]
每个 $C_n(S)$ 是复形，且
\[
C^{\cdot}\Bigl(\fU,\bigoplus_{n=-\infty}^{\infty}\cO_X(n)\Bigr)=\dlim_n C_n(S).
\]
所以我们有
\[
\bigoplus_{n=-\infty}^{\infty}H^p(X,\cO_X(n))\cong\dlim_n H^p(C_n^{\cdot}(S)).
\]
注意对每个 $p\ge 1$，序列
\[
C_n^{p-1}(S)\xrightarrow{\;d\;}C_n^p(S)\xrightarrow{\;d\;}C_n^{p+1}(S)
\]
可以等同于序列
\[
K^p(x_0^n,\dots,x_r^n;S)\xrightarrow{\;d\;}K^{p+1}(x_0^n,\dots,x_r^n;S)
\xrightarrow{\;d\;}K^{p+2}(x_0^n,\dots,x_r^n;S).
\]
于是由 \ref{it:2.4.3}，有
\[
H^p(C_n^{\cdot}(S))=H^{p+1}(x_0^n,\dots,x_r^n;S)=0
\]
当 $p\ne 0,r$ 时。因此 $\bigoplus_{n=-\infty}^{\infty}H^p(X,\cO_X(n))=0$ 对每个 $p\ne 0,r$ 成立。
\end{proof}

\begin{theorem}[Serre]\label{it:2.4.12}
设 $A$ 是诺特环，$X$ 是 $\Spec A$ 上的射影概形，$\cO_X(1)$ 是 $\Spec A$ 上丰沛的可逆 $\cO_X$-模，$\cF$ 是凝聚 $\cO_X$-模。
\begin{enumerate}[label=(\roman*),leftmargin=2.2em]
  \item （有限性定理）对每个 $p$，$H^p(X,\cF)$ 是有限生成 $A$-模。
  \item （消没定理）存在依赖于 $\cF$ 的整数 $n(\cF)$，使得对所有 $p\ge 1$ 与所有 $n\ge n(\cF)$，有 $H^p(X,\cF(n))=0$，其中 $\cF(n)=\cF\tensor_{\cO_X}\cO_X(1)^{\tensor n}$。
\end{enumerate}
\end{theorem}

\begin{proof}
可以找到闭浸入
\[
i\colon X\to\bP^r_A=\Proj A[x_0,\dots,x_r],
\]
使得 $i^{*}\cO_{\bP^r_A}(1)\cong\cO_X(1)$。可以验证 $i_*$ 是正合函子，并且把 $X$ 上的内射层映为 $\bP^r_A$ 上的内射层。所以对每个 $p$，有
\[
H^p(X,\cF)\cong H^p(\bP^r_A,i_*\cF).
\]
此外，由 \ref{it:1.4.26} 有
\[
i_*(\cF(n))\cong(i_*\cF)(n).
\]
因此可设 $X=\bP^r_A$。我们对 $p$ 作降序归纳来证明本定理。由 \ref{it:1.4.25}，可以找到满态射 $\cE\to\cF$，使 $\cE$ 是某个 $n$ 下有限多个 $\cO_X(n)$ 的直和。设 $\mathcal{R}$ 是这个满态射的核。则有正合序列
\[
0\to\mathcal{R}\to\cE\to\cF\to 0.
\]
由于 $X$ 可以被 $r+1$ 个仿射开子概形 $D_+(x_i)$（$i=0,\dots,r$）覆盖，由 \ref{it:2.3.2} 与 \ref{it:2.4.5}，有 $H^p(X,\cF(n))=0$ 对所有 $p\ge r+1$ 与所有 $n$ 成立。特别地，$H^p(X,\cF)$ 是有限生成 $A$-模，且当 $p\ge r+1$ 时 $H^p(X,\cF(n))=0$ 对所有 $n$ 成立。设 $p\ge 0$，并设对每个凝聚 $\cO_X$-模 $\cF$，我们已证 $H^{p+1}(X,\cF)$ 是有限生成 $A$-模，且存在依赖于 $\cF$ 与 $p+1$ 的整数 $n(\cF,p+1)$，使得 $H^{p+1}(X,\cF(n))=0$ 对所有 $n\ge n(\cF,p+1)$ 成立。特别地，$H^{p+1}(X,\mathcal{R})$ 是有限生成 $A$-模，且 $H^{p+1}(X,\mathcal{R}(n))=0$ 对所有 $n\ge n(\mathcal{R},p+1)$ 成立。由 \ref{it:2.4.11}，$H^p(X,\cE)$ 是有限生成 $A$-模。我们有正合序列
\[
\cdots\to H^p(X,\cE)\to H^p(X,\cF)\to H^{p+1}(X,\mathcal{R})\to\cdots.
\]
这蕴含 $H^p(X,\cF)$ 是有限生成 $A$-模。由 \ref{it:2.4.11}，对每个 $p\ge 1$，存在整数 $n(\cE,p)$，使得 $H^p(X,\cE(n))=0$ 对所有 $n\ge n(\cE,p)$ 成立。上同调群的长正合序列于是表明 $H^p(X,\cF(n))=0$ 对所有 $n\ge\max\bigl(n(\cE,p),n(\mathcal{R},p+1)\bigr)$ 成立。这说明对每个 $p\ge 1$，存在整数 $n(\cF,p)$，使得 $H^p(X,\cF(n))=0$ 对所有 $n\ge n(\cF,p)$ 成立。由于 $H^p(X,\cF(n))\ne0$ 只在 $0\le p\le r$ 时可能发生，若取 $n(\cF)=\max\limits_{1\le p\le r}n(\cF,p)$，则 $H^p(X,\cF(n))=0$ 对所有 $n\ge n(\cF)$ 与所有 $p\ge 1$ 成立。
\end{proof}

\begin{corollary}\label{it:2.4.13}
设 $f\colon X\to Y$ 是诺特概形的射影态射，$\cO_X(1)$ 是 $Y$ 上丰沛的可逆 $\cO_X$-模，$\cF$ 是凝聚 $\cO_X$-模。
\begin{enumerate}[label=(\roman*),leftmargin=2.2em]
  \item 对每个 $p$，$R^pf_*\cF$ 是凝聚 $\cO_Y$-模。
  \item 存在整数 $N$，使得对每个 $n\ge N$，典范态射 $f^{*}f_*\cF(n)\to\cF(n)$ 是满射。
  \item 存在整数 $N$，使得 $R^pf_*(\cF(n))=0$ 对每个 $n\ge N$ 与每个 $p\ge 1$ 成立。
\end{enumerate}
\end{corollary}

\begin{proof}
(i) 由 \ref{it:2.4.9}，$R^pf_*\cF$ 是拟凝聚的，且对 $Y$ 的每个仿射开子概形 $V=\Spec A$，有
\[
(R^pf_*\cF)|_V\cong H^p(f^{-1}(V),\cF)^{\sim}.
\]
由 \ref{it:2.4.12}(i)，$H^p(f^{-1}(V),\cF)$ 是有限生成 $A$-模。所以 $R^pf_*\cF$ 是凝聚的。

(ii) 对 $Y$ 的每个仿射开子概形 $V=\Spec A$，由 \ref{it:1.4.25}，存在整数 $N$，使得 $\cF(n)|_{f^{-1}(V)}$ 对每个 $n\ge N$ 都由 $\cF(n)(f^{-1}(V))$ 中的截面生成。$\cF(n)(f^{-1}(V))$ 中的每个截面 $s$ 都可以看作 $(f_*\cF(n))(V)$ 中的截面。用 $s'$ 记它在典范同态
\[
(f_*\cF(n))(V)\to\bigl(f_*f^{*}(f_*\cF(n))\bigr)(V)=f^{*}(f_*\cF(n))(f^{-1}(V))
\]
下的像。可以验证，$s'\in f^{*}f_*\cF(n)(f^{-1}(V))$ 在典范态射 $f^{*}f_*\cF(n)\to\cF(n)$ 下的像是原来的截面 $s\in\cF(n)(f^{-1}(V))$。所以
\[
(f^{*}f_*\cF(n))|_{f^{-1}(V)}\to\cF(n)|_{f^{-1}(V)}
\]
对每个 $n\ge N$ 是满射。用有限多个仿射开子概形覆盖 $Y$。我们看到，存在整数 $N$，使得 $f^{*}f_*\cF(n)\to\cF(n)$ 对每个 $n\ge N$ 是满射。

(iii) 对 $Y$ 的任一仿射开子概形 $V=\Spec A$，由 \ref{it:2.4.12}(ii)，存在整数 $N$，使得对每个 $n\ge N$ 与每个 $p\ge 1$，都有 $(R^pf_*\cF(n))|_{f^{-1}(V)}=0$。用有限多个仿射开子概形覆盖 $Y$。我们看到，存在整数 $N$，使得 $R^pf_*(\cF(n))=0$ 对每个 $n\ge N$ 与每个 $p\ge 1$ 成立。
\end{proof}

\begin{exercises}
\item \label{ex:2.4.1} 设 $f_1,\dots,f_r\in A$，$M$ 是 $A$-模，$1\le d\le r$。假设对每个 $1\le k\le d$，同态
\[
M\Big/\sum_{i=1}^{k-1}f_iM\to M\Big/\sum_{i=1}^{k-1}f_iM,\quad z\mapsto f_kz
\]
是单射。则 $H_p(f_1,\dots,f_r;M)=0$ 对每个 $p>r-d$ 成立，且 $H^p(f_1,\dots,f_r;M)=0$ 对每个 $p<d$ 成立。

\item \label{ex:2.4.2} 设 $X=\Spec A$ 是仿射概形，$\fa$ 是由 $f_1,\dots,f_r\in A$ 生成的理想，$Y=V(\fa)$，$U=X-Y$。
\begin{enumerate}[label=(\arabic*),leftmargin=2.4em]
  \item 对每个 $A$-模 $M$ 与每个 $p$，群 $H^p_Y(X,M^{\sim})$ 同构于复形
  \[
  0\to M\to\prod_i M_{f_i}\to\prod_{i<j}M_{f_if_j}\to\cdots\to M_{f_1\cdots f_r}\to 0
  \]
  的第 $p$ 个上同调，其中 $M$ 位于次数 $0$ 处，而复形的其余部分可以等同于 $M^{\sim}|_U$ 关于开覆盖 $\{D(f_1),\dots,D(f_r)\}$ 的 \v{C}ech 复形。
  \item 上述复形同构于 $\dlim_n K^{\cdot}(f_1^n,\dots,f_r^n)$，其中对任意 $n'<n$，过渡态射
  \[
  K^{\cdot}(f_1^{n'},\dots,f_r^{n'})\to K^{\cdot}(f_1^n,\dots,f_r^n)
  \]
  由
  \[
  (x_{i_1\dots i_p})\mapsto\bigl((f_{i_1}\cdots f_{i_p})^{n-n'}x_{i_1\dots i_p}\bigr)
  \]
  给出。
\end{enumerate}

\item \label{ex:2.4.3} 设 $R=\bC[t_1,\dots,t_n]$ 是复多项式环，$f\in R$，
\[
\Omega^k=\Bigl\{\sum_{i_1<\cdots<i_k}f_{i_1\dots i_k}\mathrm{d}t_{i_1}\cdots\mathrm{d}t_{i_k}:f_{i_1\dots i_k}\in R\Bigr\}
\]
是微分形式空间，$\mathrm{d}\colon\Omega^k\to\Omega^{k+1}$ 是外微分
\[
\mathrm{d}(f_{i_1\dots i_k}\mathrm{d}t_{i_1}\cdots\mathrm{d}t_{i_k})
=\sum_{i=1}^{n}\frac{\partial f_{i_1\dots i_k}}{\partial x_i}\mathrm{d}t_i\mathrm{d}t_{i_1}\cdots\mathrm{d}t_{i_k}.
\]
\begin{enumerate}[label=(\arabic*),leftmargin=2.4em]
  \item 证明复形
  \[
  \cdots\to\Omega^k\xrightarrow{\mathrm{d}f\wedge}\Omega^{k+1}\to\cdots
  \]
  同构于 Koszul 复形 $K^{\cdot}\bigl(\frac{\partial f}{\partial t_1},\dots,\frac{\partial f}{\partial t_n};R\bigr)$。
  \item 证明：对一切 $\omega\in\Omega^{\cdot}[\theta]$ 有
  \[
  (\theta\mathrm{d}+\mathrm{d}f\wedge)^2(\omega)=0,
  \]
  从而 $(\Omega^{\cdot}[\theta],\theta\mathrm{d}+\mathrm{d}f\wedge)$ 是复形。
  \item 设 $F^{\cdot}$ 是 $\Omega^{\cdot}[\theta]$ 上由
  \[
  F^p=\theta^p\Omega^{\cdot}_{\mathbb G_m^n/\mathbb C}[\theta]
  \]
  定义的滤过。证明 $\theta\mathrm{d}+\mathrm{d}f\wedge$ 保持该滤过，且有同构
  \[
  F^p/F^{p+1}\cong(\Omega^{\cdot},\mathrm{d}f\wedge),
  \]
  并且有谱序列
  \[
  E_1^{pq}=H^{p+q}(\Omega^{\cdot},\mathrm{d}f\wedge)\Rightarrow H^{p+q}(\Omega^n[\theta],\theta\mathrm{d}+\mathrm{d}f\wedge).
  \]
  典范态射 $d_1^{pq}\colon E_1^{pq}\to E_1^{p+1,q}$ 可以等同于
  \[
  H^{p+q}(\Omega^{\cdot},\mathrm{d}f\wedge)\to H^{p+q+1}(\Omega^{\cdot},\mathrm{d}f\wedge),\quad[\omega]\mapsto[\mathrm{d}\omega],
  \]
  其中 $\omega\in\Omega^{p+q}$ 满足 $\mathrm{d}f\wedge\omega=0$。
  \item 令
  \[
  K^{pq}=\theta^p\Omega^{p+q},\qquad d_1=\mathrm{d}f\wedge,\qquad d_2=(-1)^p\theta\mathrm{d}.
  \]
  则 $(K^{\cdot\cdot},d_1,d_2)$ 是双复形，上述谱序列是与该双复形相伴的谱序列之一。描述该双复形的另一个谱序列。
\end{enumerate}

\begin{tnote}
(3) 中两处记号照原书印法照录：$F^p$ 的下标印作 $\mathbb G_m^n/\mathbb C$，与本题开头由 $R=\bC[t_1,\dots,t_n]$ 定义的 $\Omega^{\cdot}$ 不合；谱序列右端印作 $\Omega^n[\theta]$，而 (2) 与本节其余各处均作 $\Omega^{\cdot}[\theta]$。读者按 $\Omega^{\cdot}[\theta]$ 理解即可。[依据：原书 p.194 习题 3 开头、p.195 (2)(3)]
\end{tnote}
\end{exercises}
\endgroup
